\subsection{Caracterizaciones del funtor producto tensorial}
\label{subsection:III.7.2}

\begin{env}[7.2.1]
\label{III.7.2.1}
Sea $A$ un anillo, no necesariamente conmutativo; sea $M$
(respectivamente $N$) un $A$-módulo a izquierda
(respectivamente a derecha), y sea $P$ un $\bb Z$-módulo.
Recordemos que dar un $\bb Z$-homomorfismo
\[
  v:N\otimes_A M\longrightarrow P
\]
equivale a dar una aplicación $\bb Z$-bilineal
\[
  u:N\times M\longrightarrow P
\]
tal que
\[
  u(ta,x)=u(t,ax)
  \qquad(a\in A,\ t\in N,\ x\in M),
\]
estando relacionadas ambas aplicaciones por
\[
  v(t\otimes x)=u(t,x).
\]
Por otra parte, dar $u$ equivale a dar un
\oldpage[III]{177}
$\bb Z$-homomorfismo
\[
  x\longmapsto f_x
\]
de $M$ en $\Hom_{\bb Z}(N,P)$ tal que
\[
  f_{ax}(t)=f_x(ta)
  \qquad(a\in A,\ t\in N,\ x\in M),
\]
estando relacionadas ambas aplicaciones por $u(t,x)=f_x(t)$.
\end{env}

\begin{env}[7.2.2]
\label{III.7.2.2}
Sea
\[
  T:\cat{Ab}_A\longrightarrow\cat{Ab}
\]
un funtor covariante aditivo. Para todo $A$-módulo a izquierda
$M$ vamos a definir un homomorfismo canónico de $\bb Z$-módulos,
funtorial en $M$,
\begin{env}[7.2.2.1]
\label{III.7.2.2.1}
\[
  t_M:T(A_s)\otimes_A M\longrightarrow T(M).
  \tag{7.2.2.1}\label{III.7.2.2.1.eq}
\]
\end{env}
Por \sref{III.7.2.1}, basta definir un $\bb Z$-homomorfismo
\[
  x\longmapsto t'_M(x)
\]
de $M$ en $\Hom_{\bb Z}(T(A_s),T(M))$ tal que
\[
  t'_M(ax)(y)=t'_M(x)(ya)
  \qquad(a\in A,\ x\in M,\ y\in T(A_s)).
\]
Para ello, obsérvese que $\Hom_{\bb Z}(T(A_s),T(M))$ es
canónicamente un $A$-módulo a izquierda, en virtud de la
estructura de $A$-módulo a derecha de $T(A_s)$, con
multiplicación escalar definida por
\[
  (a\cdot w)(y)=w(ya)
\]
para $a\in A$, $w\in\Hom_{\bb Z}(T(A_s),T(M))$ e
$y\in T(A_s)$. Definamos $t'_M$ como el compuesto de los dos
homomorfismos canónicos
\[
  M
  \xrightarrow{\sim}
  \Hom_A(A_s,M)
  \xrightarrow{T}
  \Hom_{\bb Z}(T(A_s),T(M)).
\]

La segunda flecha envía $w$ a $T(w)$; la primera es el
isomorfismo canónico de $A$-módulos que envía $x$ a la aplicación
$\theta_x:A_s\to M$ definida por
\[
  \theta_x(\xi)=\xi x
  \qquad(\xi\in A,\ x\in M).
\]
Como $\theta_{ax}=\theta_x\circ h_a$, se tiene
\[
  T(\theta_{ax})
  =
  T(\theta_x\circ h_a)
  =
  T(\theta_x)\circ T(h_a),
\]
y, por consiguiente, para $y\in T(A_s)$,
\[
  T(\theta_{ax})(y)
  =
  T(\theta_x)\bigl(T(h_a)(y)\bigr)
  =
  T(\theta_x)(ya),
\]
por la definición de la estructura de $A$-módulo a derecha sobre
$T(A_s)$. Esto prueba la existencia de $t_M$.

Es inmediato que este homomorfismo es funtorial en $M$: para todo
homomorfismo $w:M\to M'$ de $A$-módulos a izquierda, el diagrama
\begin{env}[7.2.2.2]
\label{III.7.2.2.2}
\[
\xymatrix{
  T(A_s)\otimes_A M
    \ar[r]^-{t_M}\ar[d]_{1\otimes w}
  & T(M)\ar[d]^{T(w)}
  \\
  T(A_s)\otimes_A M'
    \ar[r]_-{t_{M'}}
  & T(M')
}
\tag{7.2.2.2}\label{III.7.2.2.2.eq}
\]
\end{env}
es conmutativo. La funtorialidad de
\eqref{III.7.2.2.1.eq} muestra asimismo que, cuando $A$ es
conmutativo, $t_M$ es un homomorfismo de $A$-módulos
\sref{III.7.1.2}.
\end{env}

\begin{env}[7.2.3]
\label{III.7.2.3}
Cuando $A$ es conmutativo, se puede definir más generalmente un
homomorfismo canónico de $A$-módulos
\begin{env}[7.2.3.1]
\label{III.7.2.3.1}
\[
  T(N)\otimes_A M
  \longrightarrow
  T(N\otimes_A M).
  \tag{7.2.3.1}\label{III.7.2.3.1.eq}
\]
\end{env}
para todo $A$-módulo $N$: en la construcción de
\sref{III.7.2.2}, sustitúyase $\theta_x$ por el homomorfismo de
$A$-módulos
\[
  N\longrightarrow N\otimes_A M,
  \qquad
  y\longmapsto y\otimes x.
\]
Es inmediato que este homomorfismo es funtorial tanto en $M$ como
en $N$.

\oldpage[III]{178}
En particular, si $B$ es una $A$-álgebra, no necesariamente
conmutativa, existe un homomorfismo, funtorial en $M$,
\begin{env}[7.2.3.2]
\label{III.7.2.3.2}
\[
  (T(M))_{(B)}
  =
  T(M)\otimes_A B
  \longrightarrow
  T(M\otimes_A B)
  =
  T_{(B)}(M_{(B)}).
  \tag{7.2.3.2}\label{III.7.2.3.2.eq}
\]
\end{env}
Por la funtorialidad de \eqref{III.7.2.3.1.eq} en $M$, este es un
homomorfismo de $B$-módulos.

Se tiene además un diagrama conmutativo
\begin{env}[7.2.3.3]
\label{III.7.2.3.3}
\[
\xymatrix{
  T(A)\otimes_A M
    \ar[r]^-{t_M}\ar[d]
  & T(M)\ar[d]
  \\
  T_{(B)}(B_s)\otimes_B M_{(B)}
    \ar[r]_-{t_{M_{(B)}}}
  & T_{(B)}(M_{(B)})
}
\tag{7.2.3.3}\label{III.7.2.3.3.eq}
\]
\end{env}
La flecha vertical derecha es el compuesto
\[
  T(M)
  \longrightarrow
  T(M)\otimes_A B
  \longrightarrow
  T(M\otimes_A B)
  =
  T_{(B)}(M_{(B)})
\]
del homomorfismo canónico y \eqref{III.7.2.3.2.eq}. La flecha
vertical izquierda de \eqref{III.7.2.3.3.eq} es el homomorfismo
\[
  T(A)\otimes_A M
  \longrightarrow
  T_{(B)}(B_s)\otimes_B(B\otimes_A M)
  =
  T_{(B)}(B_s)\otimes_A M,
\]
donde
\[
  T(A)\longrightarrow T_{(B)}(B_s)=T(B)
\]
es $T(\rho)$, considerando $\rho$ como homomorfismo de
$A$-módulos $A\to B_{[\rho]}$.
\end{env}

\begin{lemma}[7.2.4]
\label{III.7.2.4}
Si $T$ es un funtor covariante aditivo de $\cat{Ab}_A$ en
$\cat{Ab}$ que conmuta con las sumas directas, entonces el
homomorfismo canónico $t_L$ de \eqref{III.7.2.2.1.eq} es un
isomorfismo para todo $A$-módulo libre $L$.
\end{lemma}

\begin{proof}
Escribamos
\[
  L=\bigoplus_{\alpha\in I}L_\alpha,
\]
donde cada $L_\alpha$ es isomorfo a $A_s$. La definición de $t_M$
en \sref{III.7.2.2} muestra que
\[
  t_L=\bigoplus_{\alpha\in I}t_{L_\alpha},
\]
porque
\[
  T:
  \Hom_A(A_s,L)
  \longrightarrow
  \Hom_{\bb Z}(T(A_s),T(L))
\]
es la suma directa de las aplicaciones $\bb Z$-lineales
\[
  T_\alpha:
  \Hom_A(A_s,L_\alpha)
  \longrightarrow
  \Hom_{\bb Z}(T(A_s),T(L_\alpha)),
\]
por la hipótesis sobre $T$. Solo queda, por tanto, demostrar el
lema para $L=A_s$; pero entonces $t_L$ es precisamente el
isomorfismo canónico
\[
  T(A_s)\otimes_A A_s
  \xrightarrow{\sim}
  T(A_s),
\]
válido para todo $A$-módulo a derecha.
\end{proof}

\begin{proposition}[7.2.5]
\label{III.7.2.5}
Sea $T$ un funtor covariante aditivo de $\cat{Ab}_A$ en
$\cat{Ab}$ que conmuta con las sumas directas. Las condiciones
siguientes son equivalentes:
\begin{enumerate}
  \item[(a)] $T$ es exacto por la derecha;
  \item[(b)] el homomorfismo canónico $t_M$ de
    \eqref{III.7.2.2.1.eq} es un isomorfismo para todo $A$-módulo
    a izquierda $M$;
  \item[($\mathrm{b}'$)] $T$ es semiexacto y $t_M$ es
    sobreyectivo para todo $A$-módulo a izquierda $M$;
  \item[(c)] $T$ es isomorfo a un funtor de $M$ de la forma
    $N\otimes_A M$, donde $N$ es un $A$-módulo a derecha.
\end{enumerate}
\end{proposition}

\begin{proof}
Es claro que (b) implica (c), y que (c) implica (a). Demostremos
que (a) implica (b).
\oldpage[III]{179}
Pongamos
\[
  T'(M)=T(A_s)\otimes_A M
\]
para todo $A$-módulo a izquierda $M$. Existe una sucesión exacta
\[
  L'\longrightarrow L\longrightarrow M\longrightarrow0,
\]
donde $L$ y $L'$ son $A$-módulos libres a izquierda. Como $T$ y
$T'$ son exactos por la derecha, existe un diagrama conmutativo
\[
\xymatrix{
  T'(L')\ar[r]\ar[d]^{t_{L'}}
  & T'(L)\ar[r]\ar[d]^{t_L}
  & T'(M)\ar[r]\ar[d]^{t_M}
  & 0
  \\
  T(L')\ar[r]
  & T(L)\ar[r]
  & T(M)\ar[r]
  & 0
}
\]
cuyas filas son exactas. Por \sref{III.7.2.4}, $t_L$ y
$t_{L'}$ son isomorfismos; por el lema de los cinco, también lo es
$t_M$.

Es claro que (b) implica ($\mathrm{b}'$). Para mostrar que
($\mathrm{b}'$) implica (a), basta demostrar el lema siguiente.
\end{proof}

\begin{lemma}[7.2.5.1]
\label{III.7.2.5.1}
Sean $\cat K$ y $\cat K'$ categorías abelianas, sean $F$ y $G$
funtores covariantes aditivos de $\cat K$ en $\cat K'$, y sea
\[
  f:F\longrightarrow G
\]
una transformación natural \emph{(T, I, 1.2)} tal que
\[
  f_E:F(E)\longrightarrow G(E)
\]
es un epimorfismo para todo objeto $E$ de $\cat K$. Si $F$ es
exacto por la derecha y $G$ es semiexacto, entonces $G$ es exacto
por la derecha.
\end{lemma}

\begin{proof}
Basta mostrar que, para todo epimorfismo $v:E'\to E$ en $\cat K$,
\[
  G(v):G(E')\longrightarrow G(E)
\]
es un epimorfismo. Existe un diagrama conmutativo
\[
\xymatrix{
  F(E')\ar[r]^{F(v)}\ar[d]_{f_{E'}}
  & F(E)\ar[d]^{f_E}
  \\
  G(E')\ar[r]_{G(v)}
  & G(E)
}
\]
Aquí $F(v)$, $f_{E'}$ y $f_E$ son epimorfismos; por consiguiente,
$G(v)$ también es un epimorfismo.
\end{proof}

\begin{remark}[7.2.6]
\label{III.7.2.6}
Para todo $A$-módulo a derecha $N$, pongamos
\[
  T_N(M)=N\otimes_A M
\]
para todo $A$-módulo a izquierda $M$. Así, $T_N$ es un funtor
covariante aditivo de $\cat{Ab}_A$ en $\cat{Ab}$, es exacto por
la derecha y conmuta con las sumas directas. Si se identifica
canónicamente $T_N(A_s)$ con $N$, se comprueba inmediatamente que
el homomorfismo correspondiente \eqref{III.7.2.2.1.eq} es la
identidad. Se sigue que el $A$-módulo a derecha $N$ de
\sref{III.7.2.5}(c) está determinado salvo isomorfismo único y es
canónicamente isomorfo a $T(A_s)$.

Equivalentemente, las correspondencias
\[
  T\longmapsto T(A_s)
  \qquad\text{y}\qquad
  N\longmapsto T_N
\]
definen una equivalencia \emph{(T, I, 1.2)} entre la categoría de
los $A$-módulos a derecha y la categoría de los funtores
covariantes aditivos
\[
  \cat{Ab}_A\longrightarrow\cat{Ab}
\]
que son exactos por la derecha y conmutan con las sumas directas.
\end{remark}

\begin{proposition}[7.2.7]
\label{III.7.2.7}
Sea $A$ un anillo artiniano a izquierda cuyo cociente por su
radical $\mathfrak m$ es un cuerpo $k$, y sea $T$ un funtor
covariante aditivo de $\cat{Ab}_A$ en $\cat{Ab}$ que conmuta con
las sumas directas. Las condiciones de \sref{III.7.2.5} son
también equivalentes a la condición siguiente:
\begin{enumerate}
  \item[(d)] $T$ es semiexacto y el homomorfismo
    \[
      T(\varepsilon):T(A_s)\longrightarrow T(k)
    \]
    inducido por el homomorfismo canónico
    $\varepsilon:A_s\to k$ es sobreyectivo.
\end{enumerate}
\end{proposition}

\begin{proof}
\oldpage[III]{180}
Es claro que la condición ($\mathrm{b}'$) de
\sref{III.7.2.5} implica (d). Demostremos que (d) implica
($\mathrm{b}'$). Existe un entero $n$ tal que
$\mathfrak m^n=0$. Para todo $A$-módulo $M$, pongamos
\[
  M_h=\mathfrak m^hM.
\]
Demostraremos, por inducción descendente sobre $h$, que $t_{M_h}$
es sobreyectivo. La afirmación es clara para $h=n$. Si $h<n$,
existe una sucesión exacta
\[
  0\longrightarrow M_{h+1}
  \longrightarrow M_h
  \longrightarrow M_h/M_{h+1}
  \longrightarrow0,
\]
y la hipótesis de inducción dice que $t_{M_{h+1}}$ es
sobreyectivo. Además, $M_h/M_{h+1}$ es anulado por $\mathfrak m$
y, por tanto, es un $(A/\mathfrak m)$-módulo; en otras palabras,
es una suma directa de $A$-módulos isomorfos a $k$. Como $T$
conmuta con las sumas directas, para demostrar que
$t_{M_h/M_{h+1}}$ es sobreyectivo basta demostrar que $t_k$ lo es.
El diagrama
\[
\xymatrix{
  T(A_s)\otimes_A A_s
    \ar[r]^-{t_{A_s}}\ar[d]_{1\otimes\varepsilon}
  & T(A_s)\ar[d]^{T(\varepsilon)}
  \\
  T(A_s)\otimes_A k
    \ar[r]_-{t_k}
  & T(k)
}
\]
es conmutativo. La condición (d), junto con
\sref{III.7.2.4}, muestra entonces que $t_k$ es sobreyectivo.

Para terminar la demostración, queda mostrar que si
\[
  0\longrightarrow M'
  \longrightarrow M
  \xrightarrow{v}M''
  \longrightarrow0
\]
es una sucesión exacta de $A$-módulos para la que $t_{M'}$ y
$t_{M''}$ son sobreyectivos, entonces $t_M$ es sobreyectivo.
Existe un diagrama conmutativo
\[
\xymatrix{
  T'(M')\ar[r]\ar[d]_{t_{M'}}
  & T'(M)\ar[r]\ar[d]_{t_M}
  & T'(M'')\ar[r]\ar[d]_{t_{M''}}
  & 0\ar[d]
  \\
  T(M')\ar[r]
  & T(M)\ar[r]_{T(v)}
  & T(M'')\ar[r]
  & \Coker(T(v))
}
\]
cuyas dos filas son exactas, por la hipótesis de que $T$ es
semiexacto. La hipótesis de inducción dice que $t_{M'}$ y
$t_{M''}$ son epimorfismos, mientras que la última flecha vertical
es un monomorfismo. El lema de los cinco \emph{(M, I, 1.1)}
muestra entonces que $t_M$ es un epimorfismo.
\end{proof}
